\documentclass[11pt,a4paper]{article}

\usepackage[T1]{fontenc}
\usepackage[utf8]{inputenc}
\usepackage{lmodern}

\usepackage[a4paper,margin=1in]{geometry}
\usepackage{setspace}
\setstretch{1.12}

\setlength{\parindent}{15pt}
\setlength{\parskip}{0pt}

\usepackage{lmodern}
\usepackage{amsmath,amssymb,amsfonts}
\usepackage{booktabs}
\usepackage{longtable}
\usepackage{array}
\usepackage{tabularx}
\usepackage{enumitem}
\usepackage{hyperref}
\usepackage{setspace}
\usepackage{ragged2e}

\hypersetup{
	colorlinks=true,
	linkcolor=black,
	citecolor=black,
	urlcolor=blue,
	pdftitle={How to Verify the General Theory of Cognitive Structuring},
	pdfauthor={Kostiantyn Osmolovskyi}
}

\setlist[itemize]{topsep=3pt,itemsep=2pt,parsep=0pt}
\setlist[enumerate]{topsep=3pt,itemsep=2pt,parsep=0pt}


\title{How to Verify\\ the General Theory of Cognitive Structuring\\
	\large A Short Roadmap for External Readers}
	
\author{
	Kostiantyn Osmolovskyi\\
	\small Independent Researcher, Ukraine\\
	\small ORCID: 0009-0006-3144-7237\\
	\small \texttt{constantinosmol@gmail.com}
}

\date{}

% ---------------------------------------------------------------------------
% Part of a series: General Theory of Cognitive Structuring (GTCS)
% Autor: Kostiantyn Osmolovskyi (ORCID: 0009-0006-3144-7237)
% License: Creative Commons Attribution 4.0 International (CC BY 4.0)
% Repository: https://github.com/k-osmolovskyi/k-osmolovskyi.github.io
% DOI: https://doi.org/10.5281/zenodo.19701902
% ---------------------------------------------------------------------------

\begin{document}
	\maketitle

\begin{center}
	\small
	Technical Note No. 2026-6\\
	Version 1.0
\end{center}

\vspace{0.5em}
	
	
\section{Purpose of this note}

This note provides a short verification roadmap for external readers of the
\emph{General Theory of Cognitive Structuring}. Its purpose is not to restate the theory,
but to reduce entry friction for reviewers, editors, formal verifiers, and adjacent-field
researchers who want to evaluate the structure of the series without reconstructing the entire
verification path from scratch.

The note answers three practical questions:
\begin{itemize}[leftmargin=2em]
	\item in what order the companion documents should be read;
	\item what each document is for;
	\item what should be checked at each stage of verification.
\end{itemize}

It is intended as a procedural guide rather than as a theoretical contribution.

\section{What this roadmap is for}

The verification package associated with the series contains several companion documents with different functions. Without a short roadmap, an external reader may face an unnecessary burden: they may not know where to begin, which document answers which question, or how to distinguish terminological, dependency, and formal traceability issues.

This note provides a minimal reading order and a task-oriented route through the verification materials.

\section{Recommended verification order}

The recommended order is as follows:

\begin{enumerate}[leftmargin=2em]
	\item \textbf{Glossary of Core Terms}
	\item \textbf{Acyclicity Statement and Dependency Criteria}
	\item \textbf{Dependency Tables and Node Registry}
	\item \textbf{Technical Appendix}
	\item \textbf{Primary papers of the series}
\end{enumerate}

This order is recommended because it moves from interpretive stabilization to structural
verification and only then to the paper corpus itself.

\section{What each document is for}

\subsection{1. Glossary of Core Terms}

\textbf{Function.}  
The glossary stabilizes the reading of the main concepts used across the series.

\textbf{What it helps verify.}
\begin{itemize}[leftmargin=2em]
	\item whether the key terms are being read in an architectural and regulatory sense rather
	than through ordinary-language substitution;
	\item whether terms such as admissibility, coherence, coherence representation,
	pre-symbolic admissibility, invariant, overload memory, and manifestation are being kept distinct;
	\item whether the reader is avoiding known category errors.
\end{itemize}

\textbf{Use this document first if you need to answer:}
\begin{itemize}[leftmargin=2em]
	\item What do the main terms mean in this framework?
	\item What do they explicitly not mean?
	\item Which terms are most vulnerable to accidental misreading?
\end{itemize}

\subsection{2. Acyclicity Statement and Dependency Criteria}

\textbf{Function.}  
This document specifies the sense in which the series is claimed to be non-circular.

\textbf{What it helps verify.}
\begin{itemize}[leftmargin=2em]
	\item whether the dependency logic of the series is acyclic under explicitly stated criteria;
	\item whether apparent cycles arise only from retrospective clarification, synthesis, or textual
	cross-reference;
	\item whether logical dependency is being properly separated from citation or thematic overlap.
\end{itemize}

\textbf{Use this document if you need to answer:}
\begin{itemize}[leftmargin=2em]
	\item Is the series built on circular definitional support?
	\item What counts as a dependency edge?
	\item Why do some apparent cycles not count as logical cycles?
\end{itemize}

\subsection{3. Dependency Tables and Node Registry}

\textbf{Function.}  
This document provides the concrete registry of nodes and the paper-level dependency structure
used to support the acyclicity claim.

\textbf{What it helps verify.}
\begin{itemize}[leftmargin=2em]
	\item which concepts, operators, constructions, and results function as explicit dependency nodes;
	\item where each node first appears;
	\item where each node is primarily formalized;
	\item how papers in the series depend on one another.
\end{itemize}

\textbf{Use this document if you need to answer:}
\begin{itemize}[leftmargin=2em]
	\item What exactly depends on what?
	\item Which paper introduces which major object?
	\item Where are the main layers of the theory located?
\end{itemize}

\subsection{4. Technical Appendix}

\textbf{Function.}  
This appendix provides a formal traceability layer for notation, assumptions, definitions,
operators, constructions, and result-level dependencies.

\textbf{What it helps verify.}
\begin{itemize}[leftmargin=2em]
	\item whether notation is being used consistently across the series;
	\item whether assumptions are being kept distinct from definitions and results;
	\item whether key propositions and theorems depend on explicitly available prior objects;
	\item whether cross-paper formal linkage is readable and internally controlled.
\end{itemize}

\textbf{Use this document if you need to answer:}
\begin{itemize}[leftmargin=2em]
	\item What are the formal prerequisites of a given definition or operator?
	\item Which assumptions are global and which are local to specific results?
	\item Which claims are definitional, and which are derived?
\end{itemize}

\subsection{5. Primary papers of the series}

\textbf{Function.}  
The papers remain the primary source of definitions, arguments, and formal development.

\textbf{What they help verify.}
\begin{itemize}[leftmargin=2em]
	\item the exact statement of definitions, constructions, propositions, and theorem-level results;
	\item the intended scope and limits of each paper;
	\item the specific formal or conceptual contribution of each paper to the series.
\end{itemize}

\textbf{Use the papers after the companion documents if you need to answer:}
\begin{itemize}[leftmargin=2em]
	\item Is the original claim actually stated as reconstructed?
	\item Is the local scope of the paper narrower or broader than the registry suggests?
	\item Does the paper support the dependency role assigned to it?
\end{itemize}

\section{Minimal verification pathway}

A minimal external verification of the series may proceed in four steps.

\subsection*{Step 1: Stabilize terminology}

Read the glossary first and make sure that the central terms are understood in their intended
architectural sense.

At the end of this step, the verifier should be able to distinguish at least:
\begin{itemize}[leftmargin=2em]
	\item coherence from coherence representation;
	\item admissibility from accessibility;
	\item pre-symbolic admissibility from structural admissibility of updating;
	\item overload from overload memory;
	\item manifestation from underlying architecture.
\end{itemize}

\subsection*{Step 2: Check non-circularity criteria}

Read the acyclicity statement and identify the dependency standard being used.

At the end of this step, the verifier should be able to distinguish:
\begin{itemize}[leftmargin=2em]
	\item logical dependency from textual cross-reference;
	\item refinement from prerequisite status;
	\item synthetic integration from foundational generation.
\end{itemize}

\subsection*{Step 3: Check structural dependency}

Use the dependency tables and node registry to inspect where each major concept and paper sits
within the unfolding of the series.

At the end of this step, the verifier should be able to answer:
\begin{itemize}[leftmargin=2em]
	\item what the major dependency nodes are;
	\item where they first appear;
	\item which papers function as entry points, extensions, or synthetic convergence nodes.
\end{itemize}

\subsection*{Step 4: Check formal traceability}

Use the technical appendix to verify notation, assumptions, and result-level dependency structure.

At the end of this step, the verifier should be able to determine:
\begin{itemize}[leftmargin=2em]
	\item whether major formal claims rest on explicit prior objects;
	\item whether the same symbol is used consistently across papers;
	\item whether assumption-dependent results are clearly distinguished from general architectural claims.
\end{itemize}

\section{If the verifier has limited time}

A reader working under time constraints may use the following compressed pathway:

\begin{enumerate}[leftmargin=2em]
	\item read the \textbf{Glossary of Core Terms};
	\item read the \textbf{Acyclicity Statement and Dependency Criteria};
	\item inspect the \textbf{paper-level dependency table} in the dependency companion;
	\item use the \textbf{Technical Appendix} only for notation or formal questions that remain unclear;
	\item consult the primary papers selectively, beginning with the ones directly relevant to the question under review.
\end{enumerate}

This shortened route is sufficient for many external verification tasks that concern structure,
dependency, and conceptual clarity rather than detailed proof reconstruction.

\section{If the verifier is checking formal consistency}

A reader interested primarily in formal consistency should proceed in the following order:

\begin{enumerate}[leftmargin=2em]
	\item \textbf{Technical Appendix}
	\item \textbf{Dependency Tables and Node Registry}
	\item \textbf{Acyclicity Statement and Dependency Criteria}
	\item relevant \textbf{primary papers}
	\item \textbf{Glossary}, only as needed for interpretive clarification
\end{enumerate}

This alternative route is useful when the main concern is not terminology but the relation among
definitions, operators, assumptions, and result-level statements.

\section{If the verifier is checking conceptual coherence}

A reader interested primarily in conceptual coherence should proceed in the following order:

\begin{enumerate}[leftmargin=2em]
	\item \textbf{Glossary of Core Terms}
	\item \textbf{Acyclicity Statement and Dependency Criteria}
	\item selected \textbf{primary papers}
	\item \textbf{Dependency Tables and Node Registry}
	\item \textbf{Technical Appendix}, only if a formal traceability question arises
\end{enumerate}

This route is useful when the central concern is whether the theory preserves distinctions among
its main conceptual layers.

\section{What this package does not try to do}

The present verification package is designed to support external reading and reconstruction.
It does not attempt to:
\begin{itemize}[leftmargin=2em]
	\item replace the papers themselves;
	\item provide a full proof dossier;
	\item settle all possible interpretive disputes in advance;
	\item force a single reading beyond the explicit dependency and terminology criteria;
	\item substitute for formal peer review.
\end{itemize}

Its purpose is more practical: to make external verification easier, faster, and less vulnerable
to accidental misreading.

\section{Recommended use of this roadmap}

This note may be used:
\begin{itemize}[leftmargin=2em]
	\item as the first entry point for external verifiers,
	\item as a front-page guide for a public verification package,
	\item as a procedural note accompanying the companion documents,
	\item as a short orientation text for reviewers who need to know where to begin.
\end{itemize}

\section{Closing note}

The \emph{General Theory of Cognitive Structuring} is supported by multiple companion documents because different verification tasks require different forms of support. Terminological stability, acyclic dependency structure, node-level traceability, and formal consistency are related but not identical tasks.

This roadmap is intended to make that division of labor explicit and to reduce the amount of
reconstruction that an external reader must perform before substantive verification can begin.

\end{document}